\color{unimain}{\subsection{Управление выключателем}}
\color{black}

\begin{enumerate}

\item
Функциональный блок управления выключателем предназначен для подачи сигнала на включение силового выключателя.

\item
Команда на включение выключателя \textit{«Включить»} формируется:
\begin{itemize}
    \item от команды оперативного включения выключателя (сигнал \textit{«РКВ»});
    \item от схемы полуавтоматического включения;
    \item при срабатывании АПВ Э1 или Э2.
\end{itemize}

\item
Функциональная схема управления выключателем приведена на рисунке \ref{pict_uv:uv}. Параметры настройки функции управления выключателем представлены в таблице \ref{app_set:uv}.

\vspace{2mm}
\begin{figure}[H]
\centering
\includegraphics[width=1\textwidth,height=1\textheight,keepaspectratio]{\fbpath/10.01 АУВ/УВ/img/uv.pdf}
\captionof{figure}{Функциональная схема управления выключателем}\label{pict_uv:uv}
\end{figure}

\item
Оперативное включение выключателя осуществляется с контролем наличия разрешающего сигнала от
функционального блока КСН (принцип работы которого описан в пункте \ref{func:ksn} настоящего руководства по
эксплуатации). Максимальное время, в течение которого продолжается ожидание условий оперативного
включения, регулируется уставкой \textbf{«Tожид усл. опер. вкл.»}. По истечению данной выдержки времени команда
оперативного включения сбрасывается. Выдержка времени на возврат \textbf{«Тпрод опер. вкл.»} задает время
продления сигнала оперативного включения.

\item
При необходимости может быть введен подхват команды на включение от токового реле,
установленного в цепи ЭМВ, в этом случае сигнал на включение удерживается в сработанном состоянии до
исчезновения тока через ЭМВ (пока блок-контакты выключателя не разорвут цепь включения). Ввод
осуществляется программным ключом \textbf{«Подхват вкл. от РТ ЭМВ»}.

\item
Предусмотрена возможность блокировки многократных включений на КЗ. Если при наличии сигнала
включения будет протекать ток по электромагнитам отключения (что соответствует включению на КЗ), то
сигнал \textit{«Включить»} заблокируется, пока не пропадут сигналы включения (от АПВ или оперативного
включения). Блокировка вводится программным ключом \textbf{«Блок. от многокр. вкл. на КЗ»}.

\item
При наличии сигнала \textit{«Блок. включения»}, формируемого функциональным блоком контроля силового
выключателя, включение выключателя блокируется.

\end{enumerate}